% !TeX root = ../navier-stokes-manuscript.tex
% ============================================================
% 2.1 Vortex Stretching
% ============================================================

Axial stretching lengthens one narrow material vortex-tube segment while its material volume stays fixed, so its cross-section contracts, its rotating fluid moves inward, and stretching spins the interior fluid faster. The velocity field carrying that segment obeys the momentum equation
\[
  \underbrace{\partial_t u}_{\text{local acceleration}}
  +
  \underbrace{(u\cdot\nabla)u}_{\text{nonlinear transport}}
  =
  \underbrace{-\nabla p}_{\text{pressure}}
  +
  \underbrace{\nu\Delta u}_{\text{viscous diffusion}}.
\]
Taking the curl removes pressure and gives vorticity advection and stretching.

\subsection{Vortex Stretching}\label{sub:ThePerfectlyStretchingVortex}

Locally the segment is axisymmetric, and the strain is uniform across the material segment and aligned with its axis. If \(e\) is that axis and \(L(t)\) is the segment's length, then
\[
  S:=\frac12\bigl(\nabla u+(\nabla u)^{\mathsf T}\bigr),
  \qquad
  Se=\sigma(t)e,
  \qquad
  \sigma(t)>0,
  \qquad
  \frac{\dot L(t)}{L(t)}
  =
  e\cdot Se
  =
  \sigma(t).
\]
Its material volume \(\mathcal V\) stays fixed. With effective transverse area \(A(t):=\mathcal V/L(t)\) and effective radius \(r_{\mathrm{eff}}(t):=\sqrt{A(t)/\pi}\),
\[
  \nabla\cdot u=0,
  \qquad
  \frac{d}{dt}(A L)=0,
  \qquad
  \frac{\dot A}{A}=-\sigma(t),
  \qquad
  \frac{\dot r_{\mathrm{eff}}}{r_{\mathrm{eff}}}
  =
  -\frac{\sigma(t)}{2}.
\]
The same narrowing segment carries vorticity \(\omega:=\nabla\times u\) according to
\[
  \frac{D\omega}{Dt}
  =
  S\omega+\nu\Delta\omega,
  \qquad
  \frac{D}{Dt}
  :=
  \partial_t+u\cdot\nabla,
\]
and under the alignment \(\omega=|\omega|e\),
\[
  S\omega
  =
  \sigma(t)\omega,
  \qquad
  \frac12\frac{D|\omega|^2}{Dt}
  =
  \sigma(t)|\omega|^2
  +
  \nu\,\omega\cdot\Delta\omega.
\]
Stretching spins the rotating fluid faster, viscosity may reinforce or oppose that spin-up, and squared vorticity grows only where the full right-hand side is positive. The energy identity proved in \Cref{prop:ClassicalKineticEnergyIdentity} and the antisymmetric part of \(\nabla u\) give
\[
  \norm{u(t)}_2
  \leq
  \norm{u_0}_2
  <\infty,
  \qquad
  \left|\frac12\bigl(\nabla u-(\nabla u)^{\mathsf T}\bigr)\right|_{\mathrm F}
  =
  \frac{|\omega|}{\sqrt2}
  \leq
  |\nabla u|_{\mathrm F}.
\]
Finite kinetic energy does not prevent the rotational part of the velocity gradient from concentrating inside the narrowing segment.

\divider{\NavierStokes{} Scaling}

At a time \(t_0<T_*\), set \(\ell:=r_{\mathrm{eff}}(t_0)\). For \(\lambda>1\), define
\[
  u_\lambda(x,t)
  :=
  \lambda u(\lambda x,\lambda^2t),
  \qquad
  p_\lambda(x,t)
  :=
  \lambda^2p(\lambda x,\lambda^2t).
\]
On the corresponding rescaled interval, \(u_\lambda\) is another \NavierStokes{} solution. At the corresponding rescaled time \(\lambda^{-2}t_0\), its segment has effective radius \(\lambda^{-1}\ell\); it is not a later material state of the opening segment. Every momentum term transforms as
\[
\begin{aligned}
  \partial_tu_\lambda(x,t)
  &=
  \lambda^3(\partial_tu)(\lambda x,\lambda^2t),
  \\
  (u_\lambda\cdot\nabla)u_\lambda(x,t)
  &=
  \lambda^3\bigl((u\cdot\nabla)u\bigr)(\lambda x,\lambda^2t),
  \\
  \nabla p_\lambda(x,t)
  &=
  \lambda^3(\nabla p)(\lambda x,\lambda^2t),
  \\
  \nu\Delta u_\lambda(x,t)
  &=
  \lambda^3\nu(\Delta u)(\lambda x,\lambda^2t).
\end{aligned}
\]
Rescaling multiplies every momentum term by the same cubic factor, so neither nonlinear transport nor viscosity gains relative ground.

\divider{Critical Height}

The corresponding Sobolev norms change with their order. Let \(\Lambda:=(-\Delta)^{1/2}\). Since
\[
  \widehat{u_\lambda}(\xi,t)
  =
  \lambda^{-2}\widehat u(\lambda^{-1}\xi,\lambda^2t),
\]
a change of variables gives
\[
  \norm{u_\lambda(t)}_{\dot H^s}^2
  =
  \lambda^{2s-1}
  \norm{u(\lambda^2t)}_{\dot H^s}^2.
\]
At \(s=0\), kinetic energy loses one power of \(\lambda\). At \(s=1\), the first derivatives gain one power. At \(s=\tfrac12\), the power is zero, so the half-derivative norm has critical height
\[
  H(t)
  :=
  \frac12\norm{u(t)}_{\dot H^{1/2}}^2
  =
  \frac12\norm{\Lambda^{1/2}u(t)}_2^2.
\]
For the rescaled field,
\[
  \norm{u_\lambda(t)}_2^2
  =
  \lambda^{-1}\norm{u(\lambda^2t)}_2^2,
  \qquad
  H_\lambda(t)=H(\lambda^2t),
  \qquad
  \norm{\nabla u_\lambda(t)}_2^2
  =
  \lambda\norm{\nabla u(\lambda^2t)}_2^2.
\]
Scale neutrality does not determine whether \(H\) rises or falls in time.

\divider{Nonlinear Production versus Viscous Dissipation}

The nonlinear term has signed production
\[
  \mathcal P_{1/2}[u](t)
  :=
  -\operatorname{Re}
  \pair
    {\Lambda^{1/2}u(t)}
    {\Lambda^{1/2}\bigl((u(t)\cdot\nabla)u(t)\bigr)}_{L^2}.
\]
Pressure vanishes from the pairing by incompressibility, while viscosity is the negative dissipation term \(-\nu\norm{\Lambda^{3/2}u}_2^2\). Thus
\[
  H'(t)
  +
  \nu\norm{\Lambda^{3/2}u(t)}_2^2
  =
  \mathcal P_{1/2}[u](t).
\]
The critical height rises when production exceeds dissipation and falls when dissipation exceeds production. Under rescaling,
\[
  \mathcal P_{1/2}[u_\lambda](t)
  =
  \lambda^2\mathcal P_{1/2}[u](\lambda^2t),
  \qquad
  \nu\norm{\Lambda^{3/2}u_\lambda(t)}_2^2
  =
  \lambda^2\nu\norm{\Lambda^{3/2}u(\lambda^2t)}_2^2.
\]
Rescaling multiplies production and dissipation by the same quadratic factor. Neither gains ground, and the nonlinear transition duration remains undetermined.

\divider{The Heat Clock}

For the linear heat equation \(v_t=\nu\Delta v\),
\[
  \widehat v(\xi,t+\delta t)
  =
  e^{-\nu|\xi|^2\delta t}\widehat v(\xi,t).
\]
A scale-localized velocity structure of width \(r\) occupies Fourier frequencies with magnitude comparable to \(r^{-1}\). Its viscous change becomes order one when
\[
  \nu|\xi|^2\delta t\simeq1,
\]
which gives the linear heat comparison time
\[
  \tau_\nu(r)
  :=
  \frac{r^2}{\nu}.
\]
Halving the width quarters that time, and more generally
\[
  \tau_\nu(\lambda^{-1}r)
  =
  \lambda^{-2}\tau_\nu(r).
\]
This linear comparison does not determine the lifetime of a nonlinear structure.

\Needspace{18\baselineskip}
\divider{The Zeno Cascade}

For the dyadic widths and their heat times
\[
  r_n:=2^{-n}r_0,
  \qquad
  \tau_n:=\frac{r_n^2}{\nu},
\]
the heat times shrink geometrically:
\[
  \tau_n
  =
  \frac{r_0^2}{\nu}4^{-n},
  \qquad
  \sum_{n=0}^{\infty}\tau_n
  =
  \frac{4r_0^2}{3\nu}
  <
  \infty.
\]
This finite sum can fit inside a finite interval, but it cannot realize a history. If one \NavierStokes{} field exhibits scale-localized structures at widths
\[
  r_0>r_1>r_2>\cdots\longrightarrow0
\]
and at sampled times
\[
  t_0<t_1<t_2<\cdots\uparrow T_*<\infty,
  \qquad
  t_{n+1}-t_n\asymp\tau_n,
\]
with comparison constants independent of \(n\), then that field passes through successively smaller structures, its remaining time satisfies
\[
  T_*-t_n
  \asymp
  \sum_{k=n}^{\infty}\tau_k
  =
  \frac{4r_n^2}{3\nu}.
\]
and its time gaps collapse at the finite endpoint \(T_*\). Even along that history, the first and successive Eulerian time derivatives of the same field at one fixed spatial point remain uncontrolled.
